\chapter*{Glossary}

\begin{description}
    % --- Data Concepts ---
    \item[Sample] A single data point or observation, typically consisting of one or more features and, in supervised learning, a label (think: rows in a spreadsheet).
    \item[Feature] An individual measurable property or characteristic of the data used by a model to make predictions (think: columns in a spreadsheet).
    \item[Training set] A collection of data used to teach a machine learning model how to make predictions.
    \item[Validation set] A subset of the data used to tune hyper-parameters and monitor the model's performance during training without affecting the final evaluation.
    \item[Testing set] A separate set of data used to evaluate how well the model performs on new, unseen inputs.
    \item[Label] Either the predicted or ground truth for a given sample. 
    For classification problems, this is a category like ``cat'' or ``dog'', often represented as a binary list with each class having one entry in said list. 
    For regression problems, this is one or more numerical values.
    % --- Models and Model Types ---
    \item[Model] Model: A machine learning model consists of two main parts: a function that maps inputs to outputs (called a hypothesis function), and a way to measure how well it performs (called a loss function). Together, these components allow the model to learn patterns from data and make predictions or decisions on new inputs.
    \item[Classifier] A type of model that assigns input data to one of several predefined categories or classes.
    \item[Regressor] A type of model that assigns one or more numeric output values to a set of input data.
    % --- Parameters and Training Process ---
    \item[Model Parameter] An internal value adjusted during training that helps the model improve its predictions.
    \item[Hyper-parameter] A predefined setting that controls how the model behaves during training or evaluations, such as the learning rate or the number of layers.
    \item[Optimisation] The process of adjusting the model’s internal parameters in order to minimize error and improve performance based on a chosen loss function.
    \item[Training] The process of teaching a machine learning model by providing it with data (a training set) and adjusting its parameters through optimization to minimize errors in predictions.
    % --- Evaluation and Validation ---
    \item[Accuracy] The fraction of correct predictions made by the model out of all predictions.
    \item[Failure Rate] The number of expected number of failures over time.
    \item[Time to Failure] The duration a system operates before a critical error occurs.
    \item[Survival time] The length of time a system continues to function correctly before experiencing a failure.
    \item[Cross-Validation] A method for evaluating model performance by repeatedly dividing the dataset into parts, training the model on some parts and testing it on the remaining parts to ensure generalisation.
    % --- Mathematical Foundations ---
    \item[Distance] A numerical measure of how different or similar two pieces of data are from each other, interpreted geometrically as the space between points.
    \item[Metric Space] A set in which distances between any two points can be consistently measured.
    \item[Linear] A relationship between variables that can be described using a straight line.
    \item[Linearly Separable] A condition where a straight line (or hyperplane in higher dimensions) lie on different sides of this plane.
    \item[Kernel] A mathematical function that allows computing the inner products in a higher-dimensional space without explicitly determining the coordinates in that space, often used in kernel  machines.
    \item[Gradient] A vector that points in the direction of the steepest increase of a function and is used to guide model parameter updates during training.
    % --- Reliability and Risk ---
    \item[Robustness] The ability of a model to maintain performance when facing unusual, noisy, or slightly changed inputs.
    \item[Reliability] The consistency with which a model or system performs its intended function correctly.
    \item[Safety-Critical] Refers to applications where incorrect model behaviour could lead to harm or serious consequences.
    \item[Security] The protection of a model or system from unauthorized access, manipulation, or attacks.
    \item[Privacy] The principle that individuals’ personal information—-such as their identity, behaviour, or sensitive attributes—-should be protected and not revealed by the model, either directly or indirectly, through its predictions or outputs.
    \item[Generalisation Gap] The difference in performance between the training data and real-world data, which indicates how well the model can generalise.
    \item[Fairness] The degree to which a model’s decisions are free of unjust discrimination across different groups or individuals.
    % --- Adversarial Context ---
    \item[Attack] A deliberate strategy used to induce undesired behaviour in a piece of software (\textit{e.g.}, a machine learning model).
    \item[Defence] A method or mechanism designed to protect a model from attacks or to reduce the harm caused by attacks.
    \item[Adversarial Input] Describes inputs specifically designed to deceive or confuse a machine learning model.
    \item[Benign Input] Describes inputs that are natural, expected, and not intended to deceive the model.
\end{description}

